\section{Conclusions}
\label{sec:conclusions}

We presented a surface- and flow-conditioned physics-informed neural network
for the onset of nucleate boiling in subcooled forced convection, built by
transfer-learning a frozen surface encoder from a pool-boiling PINN and adding a
dedicated flow encoder. The principal findings are:

\begin{enumerate}
\item \textbf{Accuracy and breadth.} On 351 curated ONB points spanning
$\Dh=0.35$--\SI{9.5}{mm} and $P=0.1$--\SI{16}{MPa}, the model predicts the ONB
wall superheat with a test RMSE of \SI{1.77}{K} ($R^2=0.79$) and heat flux with
\SI{133}{kW\,m^{-2}} ($R^2=0.78$), $2.1$--$2.7\times$ and over an order of
magnitude more accurate, respectively, than five classical correlations on the
same split.

\item \textbf{Hsu as a lower bound.} Imposing the Hsu nucleation criterion as a
one-sided lower bound, rather than an equality, eliminates a systematic
low-superheat bias in the high-subcooling regime, consistent with the physics
that subcooling and convection raise the required wall superheat.

\item \textbf{Explicit reduced pressure.} The first-order
$P\!\uparrow\!\Rightarrow\!\dTonb\!\downarrow$ trend is not recoverable from
property-based nondimensionalization alone; supplying reduced pressure as a
direct input feature, together with high-pressure training anchors, restores it
(predicted--pressure correlation $+0.01\!\to\!-0.69$).

\item \textbf{Uncertainty.} A five-member deep ensemble improves accuracy
($\dTonb$ RMSE \SI{1.73}{K}) and provides epistemic uncertainty; an honest
calibration analysis shows aleatoric scatter dominates and must be modelled for
fully calibrated intervals.
\end{enumerate}

The framework offers an accurate, physically consistent and uncertainty-aware
ONB predictor across an unprecedented parameter range, and demonstrates that a
pool-boiling surface representation transfers usefully to forced convection.
Future work will add a heteroscedastic variance head for calibrated aleatoric
uncertainty, broaden refrigerant and high-pressure coverage, solve the inverse
problem (inferring the active-cavity spectrum from measured incipience), and
extend toward transient and critical-heat-flux operation (Phase~3).
